\chapter*{Chapitre 00 -- Analyse dimensionnelle et remise a niveau}
\addcontentsline{toc}{chapter}{Chapitre 00 -- Analyse dimensionnelle et remise a niveau}

\section*{Objectif du chapitre}

Ce chapitre sert de remise a niveau avant les calculs CVC. Il regroupe 50 exercices courts,
organises par themes, pour consolider :

\begin{itemize}[leftmargin=1.4em, itemsep=0.3em]
  \item l'homogeneite dimensionnelle ;
  \item les conversions d'unites ;
  \item l'isolement d'inconnues ;
  \item le controle d'ordre de grandeur ;
  \item l'exploitation de tableaux et d'abaques.
\end{itemize}

\section*{Formulaire essentiel}

\subsection*{A. Grandeurs thermiques et energetiques}
\small
\renewcommand{\arraystretch}{1.12}
\begin{tabularx}{\textwidth}{|l|X|X|}
\hline
\textbf{Symbole} & \textbf{Grandeur} & \textbf{Unites utiles} \\
\hline
$\dot{Q}$ & Puissance thermique & W, kW, MW, kcal/h \\
\hline
$Q$ & Energie thermique & J, kJ, kWh, cal, kcal \\
\hline
$\Delta T$, $T$ & Temperature, ecart de temperature & K, $^\circ$C, $^\circ$F \\
\hline
$c_p$ & Chaleur specifique massique & J/(kg.K), kJ/(kg.K), cal/(g.$^\circ$C) \\
\hline
$h$ & Coefficient d'echange & W/(m$^2$.K), kcal/(h.m$^2$.$^\circ$C) \\
\hline
$A$ & Surface & m$^2$, cm$^2$, mm$^2$ \\
\hline
$\dot{m}$ & Debit massique & kg/s, g/s, kg/h, t/h \\
\hline
$\eta$ & Rendement & sans dimension, \% \\
\hline
\end{tabularx}

\normalsize
\vspace{0.8em}
\subsection*{B. Grandeurs hydrauliques et de debit}
\small
\renewcommand{\arraystretch}{1.12}
\begin{tabularx}{\textwidth}{|l|X|X|}
\hline
\textbf{Symbole} & \textbf{Grandeur} & \textbf{Unites utiles} \\
\hline
$\dot{V}$ & Debit volumique & m$^3$/s, m$^3$/h, L/min, L/s \\
\hline
$\Delta P$, $P$ & Pression, chute de pression & Pa, kPa, bar, mH$_2$O, atm \\
\hline
$\rho$ & Masse volumique & kg/m$^3$ \\
\hline
$v$ & Vitesse du fluide & m/s, km/h \\
\hline
$d$ & Diametre ou longueur caracteristique & m, cm, mm \\
\hline
$\ell$ & Longueur de conduite & m, cm, km \\
\hline
$\lambda$ & Coefficient de friction & sans dimension \\
\hline
$\mu$ & Viscosite dynamique & Pa.s, cP \\
\hline
$\nu$ & Viscosite cinematique & m$^2$/s, cSt \\
\hline
\end{tabularx}

\normalsize
\vspace{0.8em}
\subsection*{C. Grandeurs aero et constantes}
\small
\renewcommand{\arraystretch}{1.12}
\begin{tabularx}{\textwidth}{|l|X|X|}
\hline
\textbf{Symbole} & \textbf{Grandeur} & \textbf{Unites utiles} \\
\hline
$\mathrm{Re}$ & Nombre de Reynolds & sans dimension \\
\hline
$c$ & Vitesse du son & m/s, km/h \\
\hline
$\gamma$ & Indice adiabatique & sans dimension \\
\hline
$R$ & Constante des gaz & J/(kg.K) \\
\hline
$q_{\ell}$ & Flux lineique & W/m, kW/km, W/(100 m) \\
\hline
$q_s$ & Densite de flux & W/m$^2$, kW/m$^2$, W/cm$^2$ \\
\hline
$h_v$ & Humidite specifique & g/kg \\
\hline
$\mathrm{NUT}$ & Nombre d'unites de transfert & sans dimension \\
\hline
$\varepsilon$ & Efficacite d'echangeur & sans dimension, \% \\
\hline
$\mathrm{St}$ & Nombre de Strouhal & sans dimension \\
\hline
$f$ & Frequence & Hz, tr/min \\
\hline
$F$, $g$, $\pi$ & Force, gravite, constante pi & N ; m/s$^2$ ; $\approx 3{,}14159$ \\
\hline
\end{tabularx}
\normalsize

\subsection*{Repere important}
\begin{itemize}[leftmargin=1.4em, itemsep=0.3em]
  \item Le symbole $P$ peut designer une puissance (W) ou une pression (Pa).
  \item Toujours valider le sens par le contexte et les unites.
\end{itemize}

\section*{Enonces}

\subsection*{Theme A -- Verification d'homogeneite (6 exercices)}
\begin{enumerate}[label=\textbf{A\arabic*}, leftmargin=1.8cm, itemsep=0.4em]
  \item Perte de charge : $\Delta P = \lambda \dfrac{\ell}{d} \dfrac{\rho v^2}{2}$. Verifier l'homogeneite.
  \item Bilan thermique : $\dot{Q} = h A \Delta T$. Verifier l'homogeneite.
  \item Vitesse du son : $c = \sqrt{\dfrac{\gamma P}{\rho}}$. Verifier l'homogeneite.
  \item Puissance hydraulique : $P_{\mathrm{hyd}} = \dot{V} \times \Delta P$. Verifier l'homogeneite.
  \item Nombre de Reynolds : $\mathrm{Re} = \dfrac{\rho v d}{\mu}$. Verifier qu'il est sans dimension.
  \item Energie cinetique volumique : $\rho v^2 = \Delta P$. Verifier et commenter sa signification physique.
\end{enumerate}

\subsection*{Theme B -- Conversions simples (6 exercices)}
\begin{enumerate}[label=\textbf{B\arabic*}, leftmargin=1.8cm, itemsep=0.4em]
  \item Convertir $P = 18~\mathrm{kW}$ en W, kWh/h et cal/s.
  \item Convertir $\dot{V} = 720~\mathrm{m^3/h}$ en m$^3$/s, L/min et L/s.
  \item Convertir $\Delta P = 3{,}2~\mathrm{bar}$ en Pa, mH$_2$O et kPa.
  \item Convertir $T = 75\,^\circ\mathrm{C}$ en K et en $^\circ\mathrm{F}$. Expliquer la difference entre temperature absolue et $\Delta T$.
  \item Convertir $\dot{m} = 2{,}5~\mathrm{kg/s}$ en g/s et en t/h.
  \item Convertir $E = 125~\mathrm{kWh}$ en J, MJ et cal.
\end{enumerate}

\subsection*{Theme C -- Conversions composees (6 exercices)}
\begin{enumerate}[label=\textbf{C\arabic*}, leftmargin=1.8cm, itemsep=0.4em]
  \item Eau : $\dot{V} = 2{,}4~\mathrm{m^3/h}$ avec $\rho = 1000$ kg/m$^3$. Convertir en kg/s puis en t/h.
  \item Flux lineique : $q_{\ell} = 45~\mathrm{W/m}$. Exprimer en kW/km puis en W/(100 m).
  \item Flux de 12 kW sur 8 m$^2$. Determiner la densite de flux en W/m$^2$, kW/m$^2$ et W/cm$^2$.
  \item Air : $\mu = 1{,}8 \times 10^{-5}$ Pa.s et $\rho = 1{,}2$ kg/m$^3$. Calculer $\nu$ en m$^2$/s puis en cSt.
  \item Air humide : $\dot{V} = 500~\mathrm{L/min}$, $\rho = 1{,}2$ kg/m$^3$, $h_v = 8~\mathrm{g/kg}$. Determiner le debit d'air en m$^3$/h puis le debit d'eau en kg/h.
  \item Pompe : 15 kW absorbes, 12,3 kW hydrauliques. Exprimer les puissances en W et en kJ/h, puis calculer le rendement.
\end{enumerate}

\subsection*{Theme D -- Resolution d'equations lineaires (6 exercices)}
\begin{enumerate}[label=\textbf{D\arabic*}, leftmargin=1.8cm, itemsep=0.4em]
  \item $\dot{Q} = \dot{m} c_p (T_{\mathrm{in}} - T_{\mathrm{out}})$. Exprimer $T_{\mathrm{out}}$, $\dot{m}$ puis $\Delta T$.
  \item $\dot{Q} = h A (T_{\mathrm{paroi}} - T_{\mathrm{fluide}})$. Exprimer $T_{\mathrm{paroi}}$, $h$ puis $A$.
  \item Bifurcation : $\dot{V}_{\mathrm{tot}} = 100~\mathrm{L/min}$ et $\dot{V}_1 = 35~\mathrm{L/min}$. Determiner $\dot{V}_2$.
  \item $\dot{Q} = k \Delta T$. Si $\dot{Q} = 8~\mathrm{kW}$ et $\Delta T = 20~\mathrm{K}$, calculer $k$.
  \item $\dot{m} = \rho \dot{V}$. Avec $\rho = 998$ kg/m$^3$ et $\dot{V} = 0{,}05~\mathrm{m^3/s}$, calculer $\dot{m}$ puis reexprimer $\dot{V}$.
  \item $\Delta P = \lambda \dfrac{\ell}{d} \dfrac{\rho v^2}{2}$. Exprimer $v$ puis $\lambda$.
\end{enumerate}

\subsection*{Theme E -- Equations non lineaires (6 exercices)}
\begin{enumerate}[label=\textbf{E\arabic*}, leftmargin=1.8cm, itemsep=0.4em]
  \item $\dot{V} = v \pi \dfrac{d^2}{4}$. Exprimer $d$. Puis calculer avec $v = 2~\mathrm{m/s}$ et $\dot{V} = 0{,}05~\mathrm{m^3/s}$.
  \item $c = \sqrt{\gamma R T}$. Exprimer $T$ et $\gamma$. Verifier pour l'air avec $\gamma = 1{,}4$, $R = 287$ et $c \approx 340$ m/s.
  \item $\mathrm{Re} = \dfrac{\rho v d}{\mu}$. Exprimer $v$ puis $d$. Application : $\mathrm{Re} = 50000$, $\rho = 1000$, $\mu = 10^{-3}$ et $d = 0{,}05$ m.
  \item $P = \dot{V} \Delta P = v A \Delta P$. Exprimer $v$ en fonction de $P$, $A$ et $\Delta P$.
  \item Section annulaire : $\dot{V} = v \pi (R_e^2 - R_i^2)$. Exprimer $R_e$.
  \item $\varepsilon = 1 - \exp(-\mathrm{NUT})$. Exprimer $\mathrm{NUT}$ puis calculer pour $\varepsilon = 0{,}7$.
\end{enumerate}

\subsection*{Theme F -- Debits massique et volumique (6 exercices)}
\begin{enumerate}[label=\textbf{F\arabic*}, leftmargin=1.8cm, itemsep=0.4em]
  \item Pompe : $Q = 2{,}1~\mathrm{m^3/h}$. Convertir en L/min, m$^3$/s et L/s.
  \item Extracteur : $\dot{V} = 780~\mathrm{m^3/h}$ d'air avec $\rho = 1{,}2$ kg/m$^3$. Convertir en m$^3$/s, L/min, kg/s et t/h.
  \item Compresseur : 45 m$^3$/h a 1 atm. Convertir en L/min puis calculer $\dot{m}$ avec $\rho = 1{,}2$ kg/m$^3$.
  \item Fluide frigorifique : $\dot{m} = 0{,}15~\mathrm{kg/s}$ avec $\rho = 1300$ kg/m$^3$. Exprimer $\dot{V}$ en m$^3$/h, L/min et cm$^3$/s.
  \item Air humide : $\dot{V} = 120~\mathrm{m^3/h}$, $\rho = 1{,}18$ kg/m$^3$, $h_v = 10~\mathrm{g/kg}$. Determiner le debit sec et le debit d'eau en kg/h.
  \item Air rechauffe de 20~$^\circ$C a 60~$^\circ$C : $\rho_1 = 1{,}204$, $\rho_2 = 1{,}067$, debit initial 500 m$^3$/h. Calculer le nouveau $\dot{V}_2$ et commenter l'evolution de $\dot{m}$.
\end{enumerate}

\subsection*{Theme G -- Analyse dimensionnelle par exposants (6 exercices)}
\begin{enumerate}[label=\textbf{G\arabic*}, leftmargin=1.8cm, itemsep=0.4em]
  \item $\Delta P = k \rho^\alpha v^\beta \ell^\gamma d^\delta$. Determiner $\alpha$, $\beta$, $\gamma$ et $\delta$.
  \item $F = k \rho^\alpha v^\beta A^\gamma$ pour une trainee. Determiner les exposants.
  \item $\dot{V} = k A^\alpha \sqrt{g^\beta h^\gamma}$ pour un orifice. Determiner $\alpha$, $\beta$ et $\gamma$.
  \item $P = k \rho^\alpha g^\beta h^\gamma \dot{V}^\delta$ pour une pompe. Determiner les exposants.
  \item Verifier que $[\nu] = [\mathrm{m^2/s}]$ a partir de $\nu = \mu / \rho$.
  \item Verifier que le nombre de Strouhal $\mathrm{St} = \dfrac{f d}{v}$ est sans dimension.
\end{enumerate}

\subsection*{Theme H -- Ordres de grandeur (5 exercices)}
\begin{enumerate}[label=\textbf{H\arabic*}, leftmargin=1.8cm, itemsep=0.4em]
  \item Classer ces puissances : 10 W, 1 500 W, 3 kW, 100 kW. Associer chaque ordre de grandeur a un equipement plausible.
  \item Une vitesse d'air de 15 m/s dans une gaine de confort est-elle plausible ? Commenter au regard des pratiques usuelles.
  \item ECS : 5 m$^3$/h pour 20 personnes. Calculer le volume par personne et juger la valeur pour des douches de 10 min.
  \item Echangeur : $\Delta T = -3$ K. Est-ce possible ? Dans quel sens faut-il l'interpreter ?
  \item Paroi interieure a $-15\,^\circ$C alors que l'exterieur est a 5~$^\circ$C. Est-ce plausible ? Quelles consequences physiques ?
\end{enumerate}

\subsection*{Theme I -- Interpolation et abaques (5 exercices)}
\begin{enumerate}[label=\textbf{I\arabic*}, leftmargin=1.8cm, itemsep=0.4em]
  \item Tableau eau ($\rho$, $c_p$, $\mu$) a 0, 20, 40, 60, 80 et 100~$^\circ$C. Interpoler a 30, 50 et 55~$^\circ$C.
  \item Tableau air ($\rho$, $\nu$) a 0, 10, 20 et 30~$^\circ$C. Interpoler a 25~$^\circ$C.
  \item A partir d'un tableau type Moody : $\mathrm{Re} = 1000$, 2000, 5000, 10000. Interpoler $\lambda$ pour $\mathrm{Re} = 3000$.
  \item Viscosite d'un fluide a 25, 35 et 45~$^\circ$C. Estimer la valeur a 40~$^\circ$C et commenter son evolution avec la temperature.
  \item Enthalpie en fonction de T : 130 kJ/kg a 10~$^\circ$C, 145 kJ/kg a 30~$^\circ$C. Interpoler a 20~$^\circ$C et discuter l'hypothese lineaire.
\end{enumerate}

\subsection*{Theme J -- Synthese complete (5 exercices)}
\begin{enumerate}[label=\textbf{J\arabic*}, leftmargin=1.8cm, itemsep=0.4em]
  \item \textbf{Pompe :} $Q = 1{,}2~\mathrm{m^3/h}$, $\Delta P = 25\,000~\mathrm{Pa}$. Convertir, calculer la puissance hydraulique, puis la puissance moteur pour $\eta = 70$~\%. Commenter la plausibilite.
  \item \textbf{Radiateur :} 12 kW, eau de 70~$^\circ$C a 55~$^\circ$C avec $c_p = 4{,}18~\mathrm{kJ/(kg \cdot K)}$. Determiner le debit massique puis le convertir en kg/h et en L/min avec $\rho = 985$.
  \item \textbf{Conduit :} $\dot{V} = 400~\mathrm{m^3/h}$, $d = 0{,}2$ m, $\ell = 50$ m, $\lambda = 0{,}025$, $\rho = 1{,}2$, $\mu = 1{,}8 \times 10^{-5}$. Determiner le debit converti, la vitesse, le Reynolds et la perte de charge.
  \item \textbf{Recuperation :} effluent de 30~$^\circ$C a 20~$^\circ$C, $\dot{m} = 5~\mathrm{kg/s}$, $c_p = 4{,}18~\mathrm{kJ/(kg \cdot K)}$. Calculer la puissance, le debit equivalent, l'energie sur 24 h et le cout journalier a 0,12 euro/kWh.
  \item \textbf{Frigorifique R410A :} $\dot{m} = 0{,}08~\mathrm{kg/s}$, $\rho = 850$ kg/m$^3$, $\Delta h = 180$ kJ/kg. Determiner $\dot{V}$, la puissance frigorifique, le COP pour un moteur de 15 kW et verifier la coherence dimensionnelle.
\end{enumerate}

\newpage
\section*{Corriges express}

\subsection*{Resultats cles -- Themes A a C}
\begin{itemize}[leftmargin=1.4em, itemsep=0.35em]
  \item \textbf{Theme A :} les six relations sont homogenes. Pour A6, $\rho v^2$ represente une pression dynamique, donc une energie volumique.
  \item \textbf{B1 :} 18 000 W = 18 kWh/h $\approx$ 4 300 cal/s.
  \item \textbf{B2 :} 0,2 m$^3$/s = 200 L/s = 12 000 L/min.
  \item \textbf{B3 :} 320 000 Pa = 320 kPa = 32 mH$_2$O.
  \item \textbf{B4 :} 348 K = 167~$^\circ$F.
  \item \textbf{B5 :} 2 500 g/s = 9 t/h.
  \item \textbf{B6 :} $4{,}5 \times 10^8$ J = 450 MJ.
  \item \textbf{C1 :} 0,667 kg/s = 2 400 kg/h = 2,4 t/h.
  \item \textbf{C2 :} 45 W/m = 45 kW/km = 4 500 W/(100 m).
  \item \textbf{C3 :} 1 500 W/m$^2$ = 1,5 kW/m$^2$ = 0,15 W/cm$^2$.
  \item \textbf{C4 :} $\nu = 1{,}5 \times 10^{-5}$ m$^2$/s = 1,5 cSt.
  \item \textbf{C5 :} 30 m$^3$/h d'air et environ 0,3 kg/h d'eau.
  \item \textbf{C6 :} $\eta \approx 82$~\%.
\end{itemize}

\subsection*{Methode de travail}
\begin{itemize}[leftmargin=1.4em, itemsep=0.35em]
  \item Pour les themes D et E : poser l'equation litterale avant le numerique.
  \item Pour le theme F : garder visible la relation $\dot{m} = \rho \dot{V}$.
  \item Pour le theme G : ecrire d'abord les dimensions fondamentales M, L, T.
  \item Pour le theme H : verifier la plausibilite terrain de chaque resultat.
  \item Pour les themes I et J : expliciter les hypotheses et verifier l'ordre de grandeur final.
\end{itemize}

\begin{attention}[Corriges detailles]
Les themes D a J gagnent a etre traites avec une redaction complete : formule de depart, isolement de l'inconnue, conversions intermediaires, application numerique et verification finale.
\end{attention}
